\documentclass[11pt,a4paper]{article}
\usepackage[utf8]{inputenc}
\usepackage{amsmath,amssymb,amsthm}
\usepackage{listings}
\usepackage{xcolor}
\usepackage[margin=2.5cm]{geometry}
\usepackage{hyperref}

\lstset{
  language=C++,
  basicstyle=\ttfamily\small,
  keywordstyle=\color{blue}\bfseries,
  commentstyle=\color{green!60!black},
  stringstyle=\color{red!70!black},
  numbers=left,
  numberstyle=\tiny\color{gray},
  breaklines=true,
  frame=single,
  tabsize=4,
  showstringspaces=false
}

\title{IOI 1999 -- Code (Gray Codes)}
\author{Solution}
\date{}

\begin{document}
\maketitle

\section{Problem Statement}

Generate a Gray code of $N$ bits: a cyclic sequence of all $2^N$ binary strings of
length $N$ such that consecutive strings (including the last and first) differ in
exactly one bit.

\section{Solution Approach}

\subsection{Reflected Gray Code}

The \emph{reflected} Gray code is constructed recursively:
\begin{itemize}
  \item \textbf{Base:} For $N = 1$, the code is $(0, 1)$.
  \item \textbf{Step:} Given the $(N{-}1)$-bit Gray code $G_{N-1} = (g_0, g_1, \ldots, g_{2^{N-1}-1})$,
        the $N$-bit code is:
        \[
          (0g_0,\; 0g_1,\; \ldots,\; 0g_{2^{N-1}-1},\;
           1g_{2^{N-1}-1},\; \ldots,\; 1g_1,\; 1g_0).
        \]
\end{itemize}

\subsection{Direct Formula}

The $i$-th Gray code word (0-indexed) is:
\[
  g(i) = i \oplus \lfloor i/2 \rfloor,
\]
where $\oplus$ is bitwise XOR. This yields an $O(1)$-per-word method.

\begin{proof}[Proof that consecutive words differ in exactly one bit]
Consider $g(i) = i \oplus (i \gg 1)$ and $g(i{+}1) = (i{+}1) \oplus ((i{+}1) \gg 1)$.
Let $i$ have trailing bits $0\underbrace{1\cdots1}_k$. Then $i+1$ flips these $k{+}1$ bits.
We have $i \oplus (i{+}1) = 2^{k+1} - 1$ (a mask of $k{+}1$ ones) and
$(i \gg 1) \oplus ((i{+}1) \gg 1) = 2^k - 1$ (a mask of $k$ ones).
Their XOR is $g(i) \oplus g(i{+}1) = 2^k$, which has exactly one bit set.
\end{proof}

\section{C++ Solution}

\begin{lstlisting}
#include <cstdio>
using namespace std;

int main() {
    int N;
    scanf("%d", &N);

    long long total = 1LL << N;

    for (long long i = 0; i < total; i++) {
        long long gray = i ^ (i >> 1);
        for (int bit = N - 1; bit >= 0; bit--)
            putchar('0' + ((gray >> bit) & 1));
        putchar('\n');
    }

    return 0;
}
\end{lstlisting}

\section{Complexity Analysis}

\begin{itemize}
  \item \textbf{Time complexity:} $O(2^N \cdot N)$. Each of the $2^N$ code words
        requires $O(N)$ time to print.
  \item \textbf{Space complexity:} $O(1)$ beyond the output (each word is computed
        and printed on the fly).
\end{itemize}

\end{document}
